\documentclass[11pt]{article}
\usepackage[margin=1in]{geometry}
\usepackage[T1]{fontenc}
\usepackage{lmodern}

\usepackage{amsmath,amssymb}
\usepackage{booktabs}
\usepackage{graphicx}
\usepackage{multirow}
\usepackage{multicol}
\usepackage{enumitem}
\usepackage{hyperref}

\setlength{\parindent}{0pt}
\setlength{\parskip}{0.7\baselineskip}
\emergencystretch=2em

\title{Baseline Fusion and Preliminary Temporal-Citation Reranking for LongEval-Sci}
\author{Author Name}
\date{\today}

\begin{document}
\maketitle

\begin{abstract}
LongEval-Sci evaluates scientific retrieval under collection change, where a system should be effective on the current corpus and stable as documents accumulate over time. This paper reports a baseline and ablation study for LongEval-Sci 2026 Task 1. We compare official PyTerrier BM25 and Qwen dense baselines with full-text BM25, RM3 query expansion, cross-encoder reranking, and reciprocal rank fusion (RRF). We also evaluate a first temporal-citation reranking overlay using publication-time and OpenCitations-derived signals. On snapshot-1 train, full-text BM25 is the strongest single retriever (DCTR nDCG@10 = 0.3302, MAP = 0.2853), while RRF over full-text BM25 and dense title+abstract retrieval gives the best deep recall (Recall@1000 = 0.9667). In cumulative monthly evaluation, the same RRF run is the most stable high-performing model as the simulated corpus grows from March to May. The initial temporal-citation overlays do not yet improve the strongest lexical and fusion baselines, showing that temporal and citation evidence require careful calibration rather than uniform boosting. We conclude with a protocol for future leakage-safe temporal/citation modeling and update-trigger analysis.
\end{abstract}

\section{Introduction}

Scientific search collections are not static. New papers are published, metadata changes, citation links accumulate, and terminology shifts as fields mature. A retrieval system that is tuned once may remain strong under one snapshot yet become less reliable as the collection changes. This is the core motivation of LongEval, which frames information retrieval as a longitudinal evaluation problem rather than a one-time ranking task \cite{cancellieri_longeval_2025_clean}.

We study LongEval-Sci 2026 Task 1 with two goals. First, we want a strong non-temporal retrieval foundation for scientific search. Second, we want to understand whether simple temporal and citation signals help enough to justify adding them as reranking features. The distinction matters: a feature can be conceptually plausible and still hurt ranking if it overwhelms semantic or lexical relevance.

This paper therefore narrows the claim relative to our earlier project plan. We do not claim a finished adaptive retrieval system. Instead, we report completed baseline, fusion, temporal-overlay, citation-overlay, and month-split evaluations, then identify what remains to be done before temporal-citation reranking and update-trigger policies can be treated as validated contributions.

Our completed contributions are:
\begin{itemize}[leftmargin=1.5em]
    \item We compare five base LongEval-Sci retrieval systems: official BM25 title+abstract, official Qwen dense title+abstract, full-text BM25, title+abstract BM25 with RM3, and title+abstract BM25 with cross-encoder reranking.
    \item We evaluate three run-level RRF fusion models that reuse existing first-stage outputs without rebuilding indexes.
    \item We define a reproducible snapshot-1 development protocol using whole-train evaluation and cumulative monthly splits over March, March+April, and March+April+May.
    \item We evaluate first-pass temporal and citation overlays and show that naive temporal/citation scoring is not yet competitive with the strongest non-temporal baselines.
\end{itemize}

\section{Related Work}

\subsection{Longitudinal IR and LongEval}

LongEval evaluates how retrieval systems behave when test data diverges temporally from training data \cite{cancellieri_longeval_2025_clean}. This framing is especially relevant for scientific search, where document availability, terminology, and evidence signals such as citations all evolve. We follow this direction by separating absolute effectiveness from robustness under simulated collection growth.

\subsection{Scientific Retrieval Baselines}

The LongEval-Sci baseline space includes sparse retrieval, dense retrieval, and reranking. BM25 remains a strong and efficient foundation for large text collections \cite{zobel_inverted_2006}. Dense retrieval and cross-encoder reranking have also been explored in LongEval-Sci participant systems \cite{agh2025}. In our implementation, this motivates a two-view design: title+abstract for compact lexical and dense retrieval, and full text for a stronger sparse retrieval baseline.

\subsection{Fusion, Query Expansion, and Reranking}

RRF combines ranked lists by rank position rather than raw score, making it useful for merging sparse and dense runs whose scores are not directly comparable \cite{cormack2009rrf}. RM3-style pseudo-relevance feedback expands short queries with terms estimated from initially retrieved documents, addressing vocabulary mismatch \cite{lavrenko2001relevance}. OpenWebSearch LongEval systems similarly emphasize incremental improvements over existing rankings through reformulation, clustering, reranking, and historical evidence \cite{alexander2024openwebsearch,alexander2025openwebsearch}. Our fusion layer follows this philosophy: it operates over existing run files and does not create a new index.

\subsection{Temporal and Citation Evidence}

Temporal retrieval work has shown that freshness should be balanced with relevance in a query-dependent way \cite{dai_learning_2011}. TempRALM makes a related argument for temporally aware retrieval-augmented language models, adding temporal relevance without replacing the underlying retrieval index \cite{gade_tempralm_2025}. Citation-based scientific search provides another signal family: citation links, co-citation, and citation-based ranking can help identify relevant scientific literature beyond lexical matching alone \cite{janssens2015cocites,janssens2020validation,belter2017ranking,bascur2023citationclusters}. Our experiments treat these signals as reranking-time overlays, not as new first-stage indexes.

\subsection{Dynamic Index Maintenance}

The long-term update question is connected to dynamic indexing. Incremental inverted indexing and immediate-access dynamic indexing study how text indexes can support growing collections without full rebuilds \cite{brown_fast_nodate,asadi_fast_2013,moffat_efficient_2023}. Streaming vector-index maintenance studies similar problems for dense search \cite{mohoney_incremental_2024}. These works motivate future update-trigger experiments, but this paper does not yet validate an update policy.

\section{Task, Data, and Protocol}

We use LongEval-Sci 2026 Task 1 through the local \texttt{ir-datasets-longeval} cache. The development setting is \texttt{snapshot-1} train. It contains 869,902 documents and 100 training queries. We evaluate primarily with DCTR qrels (8,772 qrel rows) and also report raw qrels (1,183 qrel rows) for sensitivity analysis. Unless otherwise stated, tables in the main text use DCTR qrels.

The official snapshot ranges are March-May 2025 for snapshot-1, June-August 2025 for snapshot-2, and September-November 2025 for snapshot-3. Because training qrels are available for snapshot-1, we use snapshot-1 for model development and construct internal month-filtered views for robustness analysis.

\begin{table}[ht]
\centering
\caption{Snapshot-1 month-filtered evaluation views. Documents are filtered by \texttt{publishedDate}; qrels are retained only when the judged document is available in the cumulative document subset.}
\label{tab:monthly-protocol}
\begin{tabular}{lrrrr}
\toprule
Split & Months & Queries & Documents & Purpose \\
\midrule
March & 3 & 74 & 173,031 & early-corpus view \\
March+April & 3,4 & 92 & 343,421 & intermediate growth view \\
March+April+May & 3,4,5 & 96 & 525,293 & full month-filtered view \\
\bottomrule
\end{tabular}
\end{table}

This monthly protocol is an internal simulation rather than an official LongEval test condition. It filters the run and qrels to documents whose \texttt{publishedDate} month is included in the cumulative split. The query count changes because queries without at least one retained judgment are excluded from that split. This means the monthly results should be interpreted as a controlled development diagnostic, not as a complete official evaluation.

\section{Models}

Table~\ref{tab:model-protocol} summarizes the main retrieval systems. The suffix \texttt{\_dctr} in report files means that the run was evaluated with DCTR qrels; it does not mean the model was trained on DCTR labels.

\begin{table}[ht]
\centering
\caption{Model families and implementation details.}
\label{tab:model-protocol}
\begin{tabular}{p{0.31\linewidth}p{0.25\linewidth}p{0.35\linewidth}}
\toprule
Model & Retrieval view & Notes \\
\midrule
TA BM25 & title+abstract & PyTerrier BM25 anchor, Terrier stopwording and Porter stemming; config \texttt{official\_pyterrier}. \\
Dense(Qwen) & title+abstract & Qwen/Qwen3-Embedding-4B dense baseline via local OpenAI-compatible embedding service. \\
FT BM25 & full text & PyTerrier BM25 over full text; no deliberate truncation in our streaming index path. \\
TA BM25+RM3 & title+abstract & BM25 first pass, RM3 rewrite, BM25 second pass; fbDocs = 3, fbTerms = 10, fbLambda = 0.6. \\
TA rerank & title+abstract & BM25 candidates reranked with the MiniLM-L-12 MS MARCO cross-encoder. \\
RRF models & run files & Fuse two or three existing runs with RRF constant $k=60$; no index rebuild. \\
\bottomrule
\end{tabular}
\end{table}

For RRF, each query-document pair receives
\[
\mathrm{RRF}(d) = \sum_i \frac{1}{k + r_i(d)},
\]
where $r_i(d)$ is the rank of document $d$ in input run $i$. We evaluate three RRF variants: BM25 title+abstract + dense title+abstract, BM25 full text + dense title+abstract, and BM25 title+abstract + BM25 full text + dense title+abstract.

\section{Temporal and Citation Overlay}
\label{sec:temporal-citation}

The temporal/citation layer is implemented as a reranking overlay over the top retrieved candidates. It does not modify BM25, Qwen, or full-text indexes. For each query, base scores are min-max normalized within the reranked head. The final score is a weighted sum of normalized base relevance plus temporal and optional citation components.

\begin{table}[ht]
\centering
\caption{Temporal and citation features used in the first-pass overlay. These definitions are implementation-level, not claims that each feature is currently beneficial.}
\label{tab:feature-definitions}
\begin{tabular}{p{0.22\linewidth}p{0.50\linewidth}p{0.20\linewidth}}
\toprule
Feature & Definition & Leakage control \\
\midrule
Age & Days from document publication date to evaluation cutoff. & Uses \texttt{publishedDate}. \\
Recency & $\exp(-\mathrm{age}/h_f)$ with freshness half-life $h_f$. & Cutoff is snapshot/month end. \\
Foundation & $1-\exp(-\mathrm{age}/h_a)$ with age half-life $h_a$. & Uses only document date. \\
Update & Same decay as recency over update age when update dates are enabled. & Disabled or falls back when unreliable. \\
Novelty & Lightweight lexical score from novelty terms and query-unseen document terms. & Uses candidate text only. \\
Temporal intent & Rule-based query label: current, foundational, evolving, or survey. & Query text only. \\
Inbound citations & Count of citation edges to the document. & Edges after cutoff removed. \\
Recent citations & Inbound citations in a recent window before cutoff. & Edges after cutoff removed. \\
Citation velocity & Recent inbound citations divided by window length. & Edges after cutoff removed. \\
Foundational citation signal & $\log(1+\mathrm{inbound})\cdot(1-\mathrm{recent\ ratio})$. & Edges after cutoff removed. \\
Emerging citation signal & $\log(1+\mathrm{recent\ inbound})\cdot\mathrm{recent\ ratio}$. & Edges after cutoff removed. \\
\bottomrule
\end{tabular}
\end{table}

Citation features are computed from the local OpenCitations-derived CSV and filtered by citation edge creation time. Candidate documents without citation metadata receive zero citation features. Self-citation indicators are used when available through the citation graph fields.

\section{Results}

\subsection{Whole-Train Effectiveness}

Table~\ref{tab:whole-train} reports the strongest DCTR whole-train results. Full-text BM25 is the best single model on nDCG@10 and MAP. The RRF run combining full-text BM25 with dense title+abstract retrieval has slightly lower nDCG@10 but higher Recall@100 and Recall@1000, making it a strong candidate-generation strategy.

\begin{table}[ht]
\centering
\caption{Snapshot-1 train DCTR results.}
\label{tab:whole-train}
\begin{tabular}{lrrrrr}
\toprule
Method & nDCG@10 & nDCG@1000 & MAP & R@100 & R@1000 \\
\midrule
FT BM25 & 0.3302 & 0.5077 & 0.2853 & 0.7394 & 0.9245 \\
TA rerank & 0.3222 & 0.4821 & 0.2777 & 0.6952 & 0.8581 \\
RRF(FT,Dense) & 0.3175 & 0.5173 & 0.2749 & 0.7806 & 0.9667 \\
RRF(TA,FT,Dense) & 0.3028 & 0.5084 & 0.2572 & 0.7870 & 0.9806 \\
TA BM25 & 0.2922 & 0.4564 & 0.2573 & 0.6836 & 0.8581 \\
Dense(Qwen) & 0.2820 & 0.4483 & 0.2378 & 0.6390 & 0.8613 \\
TA BM25+RM3 & 0.2781 & 0.4510 & 0.2402 & 0.6559 & 0.8701 \\
\bottomrule
\end{tabular}
\end{table}

The same ordering is broadly consistent under raw qrels: \texttt{custom\_lexical\_fulltext} remains first on nDCG@10 (0.3637), while \texttt{rrf\_bm25\_ft\_dense\_ta} remains close behind (0.3532) and has the strongest deep recall among the highlighted runs.

\subsection{Temporal and Citation Overlay Results}

Table~\ref{tab:temporal-results} reports the first-pass temporal and citation overlays. These results are mainly negative. The temporal overlay improves the official dense baseline over dense alone (0.3069 vs. 0.2820 nDCG@10), but it does not beat full-text BM25 or the best RRF fusion run. On sparse full-text and title+abstract BM25, the simple temporal scoring severely hurts top-rank quality. Citation features also do not recover the loss in the current configuration.

\begin{table}[ht]
\centering
\caption{Temporal and citation overlays on snapshot-1 train DCTR.}
\label{tab:temporal-results}
\begin{tabular}{lrrrrr}
\toprule
Method & nDCG@10 & nDCG@1000 & MAP & R@100 & R@1000 \\
\midrule
Dense+temporal & 0.3069 & 0.4629 & 0.2551 & 0.6400 & 0.8613 \\
TA rerank+temporal & 0.2680 & 0.4307 & 0.2354 & 0.5447 & 0.8581 \\
TA rerank+temp+citation & 0.2543 & 0.4194 & 0.2220 & 0.5451 & 0.8581 \\
TA RM3+temporal & 0.0154 & 0.1842 & 0.0144 & 0.0433 & 0.8701 \\
FT BM25+temporal & 0.0088 & 0.1935 & 0.0091 & 0.0213 & 0.9245 \\
TA BM25+temporal & 0.0053 & 0.1750 & 0.0084 & 0.0298 & 0.8581 \\
\bottomrule
\end{tabular}
\end{table}

This is an important failure mode. The temporal overlay preserves much of the tail recall because only the head is reranked and the tail remains available, but it can move nonrelevant or weakly relevant recent/foundational documents into the top ranks. The result supports a narrower conclusion: temporal/citation evidence should be calibrated, query-conditioned, and ablated carefully before it is treated as a central ranking signal.

\subsection{Cumulative Monthly Robustness}

Table~\ref{tab:monthly-results} reports selected cumulative monthly DCTR results. The best model across all three monthly splits is \texttt{rrf\_bm25\_ft\_dense\_ta}. It also has the smallest relative improvement from March+April to March+April+May among strong systems, which means it is already strong early and remains strong as the corpus grows.

\begin{table}[ht]
\centering
\caption{Cumulative monthly DCTR nDCG@10.}
\label{tab:monthly-results}
\begin{tabular}{lrrr}
\toprule
Method & March & March+April & March+April+May \\
\midrule
RRF(FT,Dense) & 0.2732 & 0.2849 & 0.2930 \\
Dense+temporal & 0.2241 & 0.2792 & 0.2917 \\
FT BM25 & 0.2314 & 0.2259 & 0.2495 \\
TA rerank & 0.1728 & 0.2172 & 0.2655 \\
TA BM25 & 0.1767 & 0.2047 & 0.2377 \\
Dense(Qwen) & 0.2009 & 0.2182 & 0.2316 \\
\bottomrule
\end{tabular}
\end{table}

For pivot-relative temporal-change analysis, we use \texttt{official\_pyterrier\_dctr} as the pivot and nDCG@10 as the score. From March to March+April+May, the pivot improves from 0.1767 to 0.2377, giving RI = -0.3449. The best RRF run improves from 0.2732 to 0.2930, giving RI = -0.0725 and DRI = 0.2724. In this internal protocol, lower absolute RI magnitude for RRF reflects a more stable already-high score, while the pivot shows larger relative movement because it starts much lower.

\section{Discussion}

The main positive result is straightforward: full-text BM25 is the strongest current single-system foundation, and RRF over full-text BM25 plus dense title+abstract retrieval is the strongest robust candidate-generation strategy. This supports a practical architecture in which expensive or experimental reranking layers operate over a strong full-text lexical backbone and a fused candidate pool.

The temporal/citation result is more cautionary. Some temporal information helps the dense baseline, but uniform temporal scoring damages sparse systems badly. This suggests that temporal and citation features should not be globally added as bonuses. A stronger next system should use at least one of three safeguards: query-intent routing, learned score calibration, or feature-specific ablation over query groups.

The monthly protocol is useful but limited. It filters documents and qrels by \texttt{publishedDate}, so it tests robustness under a controlled development view rather than the full official future-snapshot condition. It also changes query counts across splits because some queries lose all retained judgments. We therefore treat monthly results as model-selection diagnostics, not final longitudinal claims.

\section{Limitations}

This draft has several important limitations. First, temporal and citation overlays are evaluated, but router-based and additive learned integration are not yet implemented as completed experiments. Second, citation leakage is controlled by filtering citation edges by creation time, but the reliability and coverage of citation timestamps still need a deeper audit. Third, no statistical significance tests are reported yet; future revisions should use paired tests over queries for nDCG@10 and MAP. Fourth, the monthly split uses filtered snapshot-1 judgments and may be affected by judgment incompleteness. Finally, the update-trigger section remains a design direction rather than a validated policy.

\section{Future Work: Update Triggers}

The next phase is to treat retrieval maintenance as a decision problem. Candidate trigger signals include top-k overlap drift, nDCG/MAP drops on held-out temporal splits, the fraction of newly published documents, citation-graph change, and disagreement between BM25, dense, and RRF rankings. Candidate actions should range from light to heavy: do nothing, refresh RRF, rerun temporal/citation reranking, build a delta index, refresh the full-text index, refresh the dense index, or retrain a reranker. The key open question is not simply whether to reindex, but which update action is worth its cost under the observed drift.

\section{Conclusion}

We presented a LongEval-Sci baseline and ablation study over sparse retrieval, dense retrieval, query expansion, reranking, RRF fusion, and first-pass temporal/citation overlays. The strongest current single model is full-text BM25, while the strongest robust candidate-generation strategy is RRF over full-text BM25 and dense title+abstract retrieval. Preliminary temporal/citation overlays do not yet improve the strongest baselines, and in several sparse settings they sharply degrade top-rank effectiveness. The resulting lesson is useful: before temporal and citation evidence can become a central ranking contribution, it must be calibrated and evaluated under leakage-safe temporal protocols.

\bibliographystyle{plain}
\bibliography{longeval2026_yyd}

\end{document}
